
\section{Executive Summary and Scope}
A workflow follows application-defined control edges; an agent selects some actions through model-mediated decisions while the harness enforces limits. A multi-agent system is not automatically more autonomous than a single agent. The canonical workflow motifs are prompt chaining, routing, parallelization, orchestrator--workers, and evaluator--optimizer \cite{anthropic}. ReAct-style tool loops, plan-and-execute, and peer collaboration specify additional agent execution and delegation choices.

\begin{infobox}[Preferred engineering posture]
Start with a deterministic application, a structured single model call, or a bounded chain. Introduce iterative tools only when the task depends on changing external evidence. Add workers only when decomposition is demonstrably useful. Keep action authorization, budgets, approvals, and final verification outside the model.
\end{infobox}

This guide provides original, editable diagrams and a measurable design-review framework. Performance and cost rankings are qualitative design hypotheses, not externally measured benchmarks. The target readers are AI engineers, application/platform architects, security engineers, and architecture reviewers.

\subsection{Four architectural layers}
\begin{figure}[htbp]
\centering
\resizebox{0.93\linewidth}{!}{%
\begin{tikzpicture}
\node[flowwide,text width=108mm,minimum height=12mm] (top) at (0,3) {Workflow topology: chain, parallel, route, conditional and feedback edges};
\node[flowwide,text width=108mm,minimum height=12mm] (agent) at (0,1.7) {Agent behavior: tool loop, planner, supervisor, worker, evaluator};
\node[flowcontrol,text width=108mm,minimum height=12mm] (control) at (0,0.4) {Harness and application authority: identity, permissions, policy, state, budget, trace};
\node[flowwide,text width=108mm,minimum height=12mm] (wire) at (0,-0.9) {Connectivity: model providers, MCP tool endpoints, A2A peers and data services};
\draw[flowline] (top)--(agent);
\draw[flowline] (agent)--(control);
\draw[flowline] (control)--(wire);
\end{tikzpicture}}
\caption{Four architectural layers. Neither a model proposal nor a protocol connection grants action authority.}
\label{fig:taxonomy}
\end{figure}

\subsection{Proposal is not authority}
\begin{table}[htbp]\small\centering
\begin{tabularx}{\linewidth}{@{}p{29mm}Y Y@{}}\toprule
\textbf{Concern} & \textbf{Model may propose} & \textbf{Runtime or application owns}\\\midrule
Task planning & Steps, priorities, dependency guesses & Scope, valid graph, quotas, cancellation\\
Tool use & Tool and structured arguments & Schema, caller identity, per-resource authorization\\
State & Summaries and candidate memories & Authoritative store, access and retention\\
Assessment & Evidence interpretation and critique & Objective acceptance tests and documented policy\\
High-impact action & Change request and rationale & Exact approval, version recheck, execution and audit\\
Completion & Natural-language status claim & External verification and terminal outcome\\\bottomrule
\end{tabularx}
\caption{Control-plane boundary applied consistently to all patterns.}
\end{table}

\subsection{Shared execution and state contract}
Every step should have a task ID, versioned typed input/output, principal, capability scope, dependency identifiers, timeout, resource budget, retry classification, and trace context. Use explicit states such as \texttt{PENDING}, \texttt{RUNNING}, \texttt{WAITING\_APPROVAL}, \texttt{COMPLETED}, \texttt{FAILED}, \texttt{CANCELLED}, and \texttt{TIMED\_OUT}. Do not infer success from a fluent response. A workflow can call an agent; an agent can launch a workflow; neither eliminates the need for these contracts.

\subsection{Diagram notation}
Solid navy arrows denote control or invocation. Dashed arrows denote observations, artifacts or feedback. Teal boxes represent deterministic validators or governance services. Pale-red boxes represent human or elevated-risk gates. Diagrams are logical views, not deployment diagrams; individual products and hosting choices are intentionally omitted.
